\documentclass[../DoAn.tex]{subfiles}
\begin{document}

\section{Tổng quan kỹ thuật}
\label{section:5.1}
 Nhằm đảm bảo trải nghiệm chơi game đi cảnh hành động nhịp độ cao diễn ra mượt mà nhất, trò chơi hướng tới mục tiêu hiệu năng đạt mức \textbf{60 FPS} ổn định trên cấu hình khuyến nghị. Mức khung hình này giúp đồng bộ hóa tối ưu giữa chuỗi xử lý lệnh di chuyển và phản hồi hoạt ảnh của nhân vật, giảm thiểu tối đa hiện tượng trễ hình (input lag) hay giật xé hình khi người chơi thực hiện các thao tác di chuyển phức tạp.

Hệ thống yêu cầu về phần cứng máy tính tối thiểu và khuyến nghị để vận hành trò chơi được chi tiết hóa trong bảng dưới đây:

\begin{table}[H]
\centering
\begin{tabular}{|p{3.0cm}|p{4.5cm}|p{5.0cm}|}
\hline
\textbf{Thành phần phần cứng} & \textbf{Cấu hình tối thiểu (Minimum)} & \textbf{Cấu hình khuyến nghị (Recommended)} \\ \hline
\textbf{Hệ điều hành (OS)} & Windows 10/11 (64-bit) & Windows 10/11 (64-bit) \\ \hline
\textbf{Bộ vi xử lý (CPU)} & Intel Core i3-4160 hoặc AMD Phenom II X4 965 & Intel Core i5-6400 hoặc AMD Ryzen 3 1200 \\ \hline
\textbf{Bộ bộ nhớ trong (RAM)} & 4 GB RAM & 8 GB RAM \\ \hline
\textbf{Card đồ họa (GPU)} & NVIDIA GeForce GTX 560 hoặc AMD Radeon HD 5770 (1GB VRAM) & NVIDIA GeForce GTX 960 hoặc AMD Radeon R9 380 (2GB VRAM) \\ \hline
\textbf{Phiên bản DirectX} & DirectX 11 & DirectX 11 \\ \hline
\textbf{Dung lượng ổ cứng} & 1 GB dung lượng trống & 1 GB dung lượng trống (ưu tiên SSD) \\ \hline
\end{tabular}
\caption{Bảng yêu cầu cấu hình phần cứng tối thiểu và khuyến nghị}
\label{table:hardware_requirements}
\end{table}


\section{Thiết kế kiến trúc}
\label{section:5.2}

\subsection{Lựa chọn kiến trúc phần mềm}
\label{subsection:5.2.1}
Để phát triển hệ thống điều khiển nhân vật và quái vật linh hoạt trong thể loại Metroidvania hành động nhịp độ cao, đề tài lựa chọn áp dụng mô hình \textbf{Kiến trúc hướng thành phần (Component-Based Architecture)} đặc thù của Unity Engine kết hợp với \textbf{Mẫu thiết kế Trung gian (Mediator Pattern)} và \textbf{Kiến trúc dữ liệu hướng cấu hình (Data-Driven Design)} để tổ chức cấu trúc lớp trong trò chơi.

Trước khi đi sâu vào thiết kế thực thể người chơi, cần làm rõ hai khái niệm cốt lõi của Kiến trúc hướng thành phần trong Unity:
\begin{itemize}
    \item \textbf{GameObject (Đối tượng trò chơi):} Là thực thể cơ sở và là container trống đóng vai trò chứa đựng các thành phần logic. Bản thân GameObject không tự sở hữu bất kỳ hành vi hay dữ liệu nào ngoại trừ các thông số vật lý cơ bản về vị trí, góc xoay và tỉ lệ (được định nghĩa trong component mặc định \texttt{Transform}). Tất cả các thực thể trong trò chơi như nhân vật chính, kẻ địch, cạm bẫy hay camera đều được biểu diễn dưới dạng GameObject.
    \item \textbf{Component (Thành phần chức năng):} Là các khối logic mã nguồn độc lập (kế thừa từ lớp cơ sở \texttt{MonoBehaviour} của Unity) được gắn trực tiếp vào GameObject để cung cấp các thuộc tính và hành vi chuyên biệt. Mô hình này tuân thủ nguyên lý thiết kế ưu tiên lắp ghép hơn kế thừa sâu lớp (Composition over Inheritance). Bằng việc kết hợp các thành phần khác nhau, nhà phát triển có thể tạo ra các thực thể có hành vi phong phú mà không gặp phải sự cứng nhắc và phức tạp của cây kế thừa lớp truyền thống.
\end{itemize}

Áp dụng nguyên lý trên vào thực thể chính của người chơi (\texttt{GameObject} Player), hệ thống được lắp ghép từ các component hoạt động độc lập dưới đây:
\begin{itemize}
    \item \texttt{PlayerInput}: Chịu trách nhiệm đọc, lọc nhiễu và lưu vào bộ đệm các tín hiệu nhập liệu bàn phím/chuột từ người dùng (ví dụ: Jump buffering để đệm lệnh nhảy trước khi chạm đất, Dash buffering).
    \item \texttt{PlayerMovement}: Xử lý động học, tính toán lực nhảy, di chuyển, lướt 8 hướng, bám vách tường và tương tác vật lý với thế giới thông qua Rigidbody2D.
    \item \texttt{PlayerStats}: Quản lý các chỉ số sinh tồn runtime thay đổi liên tục (Máu, Thể lực, Năng lượng) và kích hoạt các sự kiện C\# event khi các chỉ số này thay đổi.
    \item \texttt{PlayerCollision}: Chịu trách nhiệm kiểm tra va chạm vật lý với môi trường địa hình, vùng sát thương của bẫy gai và nhận diện các trigger an toàn.
    \item \texttt{PlayerCombat}: Quản lý trạng thái chiến đấu, chuỗi combo 3-hit cận chiến và kích hoạt bật/tắt các hitbox tấn công động.
    \item \texttt{PlayerFormManager}: Quản lý cơ chế biến đổi linh hồn, thực hiện nạp và hủy động các cấu trúc Prefab vật lý của hình thái biến đổi.
\end{itemize}

Để các component trên giao tiếp với nhau hiệu quả mà không bị liên kết chặt (tight coupling), trò chơi áp dụng \textbf{Mẫu thiết kế Trung gian (Mediator Pattern)}:
\begin{itemize}
    \item Lớp \texttt{PlayerController} đóng vai trò là Mediator trung tâm. Trong pha khởi tạo (\texttt{Awake}), lớp này lưu trữ tham chiếu đến tất cả các component con trên.
    \item Khi có sự kiện phát sinh, các component con không gọi trực tiếp lẫn nhau mà gửi thông điệp thông qua \texttt{PlayerController}. Ví dụ, khi \texttt{PlayerInput} nhận tín hiệu nhảy, nó thông báo cho \texttt{PlayerController}, từ đó điều phối gọi phương thức nhảy trong \texttt{PlayerMovement} chỉ khi \texttt{PlayerCollision} xác nhận nhân vật đang đứng trên mặt đất.
\end{itemize}

Đồng thời, trò chơi áp dụng \textbf{Kiến trúc dữ liệu hướng cấu hình (Data-Driven Design)} thông qua việc sử dụng \textbf{ScriptableObject}:
\begin{itemize}
    \item ScriptableObject là một thùng chứa dữ liệu tĩnh đặc thù của Unity, tồn tại độc lập dưới dạng tài sản tệp tin trong bộ nhớ dự án mà không cần gắn vào bất kỳ GameObject nào.
    \item Việc áp dụng ScriptableObject (như lớp \texttt{PlayerFormDataSO} cấu hình thuộc tính linh hồn, \texttt{ItemDatabaseSO} định nghĩa danh sách vật phẩm, \texttt{SpellDatabaseSO} chứa các kỹ năng phép thuật) giúp phân tách hoàn toàn dữ liệu cân bằng trò chơi ra khỏi logic xử lý. Nhà thiết kế game có thể trực tiếp tinh chỉnh các thông số như HP tối đa, tốc độ chạy, hoặc gán các prefab vật lý khác nhau trực tiếp trên giao diện Unity Inspector mà không cần sửa đổi hay biên dịch lại mã nguồn hệ thống.
\end{itemize}


\subsection{Thiết kế tổng quan}
\label{subsection:5.2.2}
Để tổ chức mã nguồn rõ ràng và dễ bảo trì, kiến trúc phần mềm của trò chơi được phân rã thành các gói. Sơ đồ phân rã các gói chính trong mã nguồn trò chơi được trình bày tại Hình~\ref{fig:system_decomposition}.

\begin{figure}[H]
    \centering
    \includegraphics[width=0.85\textwidth]{Hinhve/system_decomposition.png}
    \caption{Sơ đồ phân rã các gói chính trong mã nguồn trò chơi}
    \label{fig:system_decomposition}
\end{figure}

\textit{Giải thích Sơ đồ phân rã các gói (Hình~\ref{fig:system_decomposition}):}
Hệ thống mã nguồn trò chơi được chia thành 10 gói chức năng chính cụ thể như sau:
\begin{itemize}
    \item \textbf{Gói Player:} Quản lý toàn bộ vòng đời, máy trạng thái di chuyển, hoạt ảnh và các hành vi vật lý/chiến đấu của nhân vật chính. Gói chứa các thành phần con xử lý di chuyển nền tảng (Platforming), lướt né tránh (Dash), nhảy đôi, bám vách leo tường và cơ chế chiến đấu cận chiến 3-hit combo. Đồng thời, gói cũng điều khiển hệ thống chuyển đổi hình thái linh hồn (Form Switch), giúp nhân vật thay đổi cấu hình Animator và Prefab vật lý theo từng dạng quái vật hấp thụ được. Bên cạnh đó, hệ thống tích hợp các hiệu ứng âm thanh (SFX) phong phú cho từng hành vi của nhân vật (như nhảy, lướt, nhận sát thương, hồi sinh) cùng nhạc nền (BGM) nhằm nâng cao tính chân thực và trải nghiệm nghe nhìn của người chơi.
    \item \textbf{Gói Enemy:} Thiết lập và điều phối hành vi của toàn bộ kẻ địch thông thường cùng các thực thể Boss lớn trong trò chơi. Phân hệ này sử dụng cấu trúc AI Brain kết hợp máy trạng thái (StateMachine) để quái vật tuần tra ngẫu nhiên, chuyển sang trạng thái cảnh giác khi nghe tiếng động và truy đuổi người chơi khi bước vào tầm quét. Đối với Boss, gói quản lý trọng số mong muốn (Desire Weights) và các chuỗi chiêu thức liên hoàn (Combo Patterns) phức tạp.
    \item \textbf{Gói System:} Quản lý các dịch vụ và hệ thống nền tảng phục vụ cho vận hành game cục bộ. Gói bao gồm phân hệ lưu trữ và tải game (Save/Load) dưới định dạng JSON bền vững trên thiết bị, hỗ trợ nhiều ô lưu trữ (Save Slot) giúp ghi lại thời gian chơi thực tế, tiến trình nâng cấp làng (Facility) và trạng thái mở khóa trên Cây kỹ năng (Skill Tree). Ngoài ra, gói còn quản lý hệ thống nâng cấp thuộc tính vĩnh viễn của nhân vật, trạng thái mua sắm tại NPC Shopkeeper (Shop) và hệ thống tối ưu hóa hiệu năng culling tắt/bật có chọn lọc các thành phần vật lý và logic dựa trên khoảng cách (Distance Culling) và vùng phòng chơi (Room Culling).
    \item \textbf{Gói UI:} Quản lý các Canvas giao diện hiển thị tĩnh và động trong trò chơi. HUD Canvas hiển thị lượng máu, thể lực, năng lượng thời gian thực của người chơi. Panels Canvas hiển thị các bảng menu chính, menu lưu trữ, túi đồ, cây kỹ năng nâng cao và bản đồ định vị thế giới. World Canvas được sử dụng để hiển thị các thành phần UI động như chữ sát thương bay (Damage Text) tại tọa độ không gian thế giới. Các thông báo phản hồi và giao diện xác nhận (Confirm Panel) được thống nhất thiết kế bằng tiếng Anh để đảm bảo tính đồng bộ và chuyên nghiệp.
    \item \textbf{Gói Interfaces:} Định nghĩa các hợp đồng giao tiếp trừu tượng phục vụ cho việc liên kết lỏng (loose coupling) giữa các phân hệ. Tiêu biểu là giao diện nhận sát thương (\texttt{IDamageable}) được thực thi bởi cả người chơi và quái vật để nhận lực đẩy lùi và giảm máu, cùng giao diện tương tác (\texttt{IInteractable}) cho phép người chơi tương tác với cửa, hòm đồ, xác quái vật và đền thờ.
    \item \textbf{Gói Core:} Đóng vai trò là nền tảng khởi tạo và điều phối dòng chạy chính của trò chơi. Gói chứa các cấu hình hằng số toàn cục, các lớp quản lý theo mẫu Singleton (như GameManager, UIManager, GamePauseManager), và các tiện ích dùng chung xuyên suốt các màn chơi để đồng nhất dữ liệu.
    \item \textbf{Gói Trap:} Quản lý các đối tượng cạm bẫy vật lý được bố trí trong môi trường bản đồ như bẫy gai nhọn, lưỡi dao xoay. Các lớp trong gói này chịu trách nhiệm kiểm tra va chạm vật lý với người chơi và kích hoạt trừ máu trực tiếp thông qua giao diện nhận sát thương, góp phần tạo nên thử thách đi cảnh.
    \item \textbf{Gói VFX:} Quản lý toàn bộ hiệu ứng hình ảnh, hoạt ảnh phản hồi và các hệ thống hạt (Particle System) trong game. Phân hệ này kích hoạt các hiệu ứng đặc biệt khi người chơi lướt né đòn, thực hiện combo chém cận chiến, thi triển đạn phép thuật hoặc khi quái vật tan biến khi bị tiêu diệt, giúp tăng trải nghiệm thị giác.
    \item \textbf{Gói Scriptable Object:} Lưu trữ và quản lý toàn bộ cấu trúc dữ liệu hướng cấu hình (Data-Driven Design) của dự án. Phân hệ này chứa các tài sản dữ liệu tĩnh như cơ sở dữ liệu vật phẩm, phép thuật, cây kỹ năng, cấu hình hình thái linh hồn (Form Switch), và các trạng thái chạy runtime dùng chung như ví tiền vàng, trạng thái hòm đồ nâng cấp.
    \item \textbf{Gói Interact Object:} Quản lý các đối tượng tương tác vật lý trong thế giới game thông qua va chạm trigger 2D. Phân hệ bao gồm các điểm nhặt vật phẩm rơi ra (Item Pickup), rương sách phép thuật (Spell Book Chest), rương lưu trữ đồ (Storage), cần gạt kích hoạt thang máy, và các vùng hiển thị gợi ý phím tương tác lên màn hình HUD của người chơi. Đồng thời, phân hệ tích hợp cơ chế xác suất rơi vật phẩm ngẫu nhiên khi tiêu diệt kẻ địch, trong đó các nguyên liệu rơi ra (như các vật phẩm nguyên liệu A, B, C, D) được thu thập để làm tài nguyên nâng cấp các chỉ số tại cơ sở nâng cấp Facility.
\end{itemize}

Tiếp theo, sơ đồ kiến trúc luồng dữ liệu biểu thị các liên kết tương tác giữa các lớp logic được trình bày tại Hình~\ref{fig:game_architecture}.

\begin{figure}[H]
    \centering
    \includegraphics[width=0.75\textwidth]{Hinhve/game_architecture.png}
    \caption{Kiến trúc tổng thể các lớp logic trong hệ thống}
    \label{fig:game_architecture}
\end{figure}

\textit{Giải thích Sơ đồ kiến trúc tổng thể các lớp logic (Hình~\ref{fig:game_architecture}):}
Luồng dữ liệu và tương tác vật lý diễn ra nhịp nhàng giữa các tầng:
\begin{itemize}
    \item Khi người chơi thao tác, thành phần \texttt{PlayerInput} thu nhận tín hiệu và gửi tới bộ điều phối trung tâm \texttt{PlayerController}.
    \item \texttt{PlayerController} đóng vai trò Mediator ra lệnh cho \texttt{PlayerMovement} áp lực vật lý, hoặc yêu cầu \texttt{PlayerFormManager} tải dữ liệu cấu hình từ các \texttt{ScriptableObjects} để biến đổi hình thái.
    \item Khi máu nhân vật thay đổi, \texttt{PlayerStats} phát đi sự kiện dạng event-driven để cập nhật lập tức lên màn hình \texttt{PlayerHudUI}.
    \item \texttt{GameSaveManager} chịu trách nhiệm thu thập trạng thái từ các component này để thực hiện tuần tự hóa, trong khi \texttt{CullingManager} phối hợp khoảng trễ (Hysteresis) và trạng thái phòng (\texttt{RoomController}) để vô hiệu hóa có chọn lọc các đối tượng nằm ngoài tầm nhìn hoặc không hoạt động nhằm tối ưu hóa hiệu năng xử lý và vẽ đồ họa.
\end{itemize}


\subsection{Thiết kế chi tiết gói}
\label{subsection:5.2.3}
Dưới đây là biểu đồ thiết kế gói chi tiết cho các gói cốt lõi của hệ thống bao gồm: gói Player (Hình~\ref{fig:player_package_design}), gói Enemy (Hình~\ref{fig:enemy_package_design}), gói UI (Hình~\ref{fig:ui_package_design}), gói System (Hình~\ref{fig:system_package_design}), gói Scriptable Object (Hình~\ref{fig:scriptable_object_package_design}), gói Interact Object (Hình~\ref{fig:interact_object_package_design}), và biểu đồ phân cấp kẻ địch (Hình~\ref{fig:enemy_class_hierarchy}).

\begin{figure}[H]
    \centering
    \includegraphics[width=0.85\textwidth]{Hinhve/player_package_design.png}
    \caption{Biểu đồ thiết kế chi tiết bên trong gói Player (Player Package Design)}
    \label{fig:player_package_design}
\end{figure}

\textit{Giải thích biểu đồ gói Player (Hình~\ref{fig:player_package_design}):}
Biểu đồ mô tả cấu trúc lắp ghép và các mối quan hệ lớp chuẩn UML bên trong phân hệ người chơi:
\begin{itemize}
    \item Lớp điều phối trung tâm \texttt{PlayerController} duy trì mối quan hệ Hợp thành (Composition) với máy trạng thái \texttt{PlayerStateMachine}.
    \item \texttt{PlayerController} liên kết thông qua quan hệ Kết hợp (Association) đến các component độc lập (\texttt{PlayerInput}, \texttt{PlayerMovement}, \texttt{PlayerStats}, \texttt{PlayerFormManager}) để phân bổ nhiệm vụ xử lý di chuyển, chỉ số và biến hình.
    \item \texttt{PlayerFormManager} sở hữu quan hệ Kết tập (Aggregation) với các đối tượng cấu hình \texttt{PlayerFormDataSO} kế thừa từ ScriptableObject.
    \item \texttt{PlayerController} thực thi (Implementation) giao diện trừu tượng \texttt{IDamageable} từ gói \texttt{Interfaces} để cho phép nhận sát thương đồng nhất.
\end{itemize}

\begin{figure}[H]
    \centering
    \includegraphics[width=0.8\textwidth]{Hinhve/enemy_package_design.png}
    \caption{Biểu đồ thiết kế chi tiết bên trong gói Enemy (Enemy Package Design)}
    \label{fig:enemy_package_design}
\end{figure}

\textit{Giải thích biểu đồ gói Enemy (Hình~\ref{fig:enemy_package_design}):}
Biểu đồ đặc tả cấu trúc máy trạng thái AI dùng chung cho toàn bộ quái vật:
\begin{itemize}
    \item Lớp cơ sở trừu tượng \texttt{MvEnemyBase} kế thừa từ \texttt{MonoBehaviour} của Unity, thực thi giao diện nhận sát thương \texttt{IDamageable} và sở hữu mối quan hệ Hợp thành (Composition) với \texttt{EnemyStateMachine}.
    \item \texttt{EnemyStateMachine} quản lý danh sách các trạng thái \texttt{EnemyState} độc lập. Mỗi \texttt{EnemyState} duy trì một tham chiếu đến \texttt{EnemyContext}.
    \item Lớp \texttt{MvEnemyBase} kết hợp (Association) với thành phần xử lý hoạt ảnh \texttt{EnemyAnimationController} và lớp tấn công vật lý \texttt{MvAttack}.
    \item Các lớp kẻ địch cụ thể như \texttt{MvEm0010} (BoneDog) sử dụng quan hệ Kế thừa (Inheritance) từ \texttt{MvEnemyBase} để tái sử dụng toàn bộ khung AI.
\end{itemize}

\begin{figure}[H]
    \centering
    \includegraphics[width=0.8\textwidth]{Hinhve/ui_package_design.png}
    \caption{Biểu đồ thiết kế chi tiết bên trong gói UI (UI Package Design)}
    \label{fig:ui_package_design}
\end{figure}

\textit{Giải thích biểu đồ gói UI (Hình~\ref{fig:ui_package_design}):}
Biểu đồ mô tả cấu trúc phân lớp Canvas và quản lý giao diện người dùng:
\begin{itemize}
    \item Lớp điều phối trung tâm \texttt{UIManager} sở hữu mối quan hệ Hợp thành (Composition) với \texttt{UIStateManager} để quản lý máy trạng thái hiển thị của UI.
    \item \texttt{UIManager} duy trì quan hệ kết hợp đến các Canvas giao diện hiển thị chuyên biệt: \texttt{PlayerHudUI} (HUD máu, thể lực), \texttt{SaveSlotMenuUI} (Menu chọn tệp lưu trữ), và \texttt{TitleCanvasManager} (Màn hình chính bắt đầu game).
    \item \texttt{SaveSlotMenuUI} liên kết với mảng các nút chọn \texttt{SaveSlotButtonView}.
    \item \texttt{TitleCanvasManager} kết hợp với thành phần thiết lập âm lượng \texttt{TitleSettingMenuController}.
    \item Lớp \texttt{DamageTextView} đại diện cho UI hiển thị chữ sát thương bay trong không gian 2D.
\end{itemize}

\begin{figure}[H]
    \centering
    \includegraphics[width=0.9\textwidth]{Hinhve/system_package_design.png}
    \caption{Biểu đồ thiết kế chi tiết bên trong gói System (System Package Design)}
    \label{fig:system_package_design}
\end{figure}

\textit{Giải thích biểu đồ gói System (Hình~\ref{fig:system_package_design}):}
Biểu đồ mô tả hệ thống quản lý lưu trữ và tiến trình của trò chơi:
\begin{itemize}
    \item Lớp quản lý lưu trữ \texttt{GameSaveManager} thiết kế theo mẫu Singleton duy trì mối quan hệ phụ thuộc với cấu trúc dữ liệu \texttt{GameSaveData} để thực hiện serialize/deserialize JSON.
    \item \texttt{GameSaveManager} kết hợp (Association) với các ScriptableObject trạng thái động như \texttt{CurrencyWalletSO} (ví tiền vàng) và \texttt{InventoryStateSO} (túi đồ), đồng thời phụ thuộc vào lớp dịch vụ tĩnh \texttt{PortalCheckpointService} để ghi nhận các điểm dịch chuyển đã mở khóa.
\end{itemize}

\begin{figure}[H]
    \centering
    \includegraphics[width=0.95\textwidth]{Hinhve/scriptable_object_package_design.png}
    \caption{Biểu đồ thiết kế chi tiết bên trong gói Scriptable Object (Scriptable Object Package Design)}
    \label{fig:scriptable_object_package_design}
\end{figure}

\textit{Giải thích biểu đồ gói Scriptable Object (Hình~\ref{fig:scriptable_object_package_design}):}
Biểu đồ mô tả cấu trúc phân lớp dữ liệu hướng cấu hình ScriptableObject trong game:
\begin{itemize}
    \item Lớp cơ sở \texttt{ScriptableObject} cung cấp các thuộc tính quản lý bộ nhớ của Unity Engine.
    \item Các lớp dữ liệu cấu hình tĩnh (như \texttt{ItemDefinitionSO} định nghĩa thuộc tính vật phẩm, \texttt{KeyItemSO} và \texttt{ToolItemSO} kế thừa để phân biệt loại vật phẩm, \texttt{SpellDatabaseSO} và \texttt{SpellData} chứa danh sách phép thuật, \texttt{SkillTreeDatabaseSO} chứa nhánh kỹ năng, \texttt{FacilityUpgradeDatabaseSO} chứa thông tin nâng cấp làng) tồn tại dưới dạng tài sản độc lập trong bộ nhớ.
    \item Các lớp lưu trữ trạng thái chạy runtime (như \texttt{InventoryStateSO} quản lý số lượng vật phẩm trong túi đồ, \texttt{CurrencyWalletSO} quản lý số vàng, \texttt{FacilityProgressSO} quản lý cấp độ nâng cấp làng hiện tại) giúp các phân hệ truy xuất nhanh mà không phụ thuộc lẫn nhau.
\end{itemize}

\begin{figure}[H]
    \centering
    \includegraphics[width=0.85\textwidth]{Hinhve/interact_object_package_design.png}
    \caption{Biểu đồ thiết kế chi tiết bên trong gói Interact Object (Interact Object Package Design)}
    \label{fig:interact_object_package_design}
\end{figure}

\textit{Giải thích biểu đồ gói Interact Object (Hình~\ref{fig:interact_object_package_design}):}
Biểu đồ mô tả cấu trúc thiết kế của phân hệ tương tác trong trò chơi:
\begin{itemize}
    \item Giao diện \texttt{IInteractable} trong gói Interfaces định nghĩa các hàm chung cho các thực thể tương tác bao gồm: \texttt{Interact} (thực thi hành động), \texttt{GetInteractPrompt} (lấy văn bản chỉ dẫn), và \texttt{CanInteract} (kiểm tra điều kiện tương tác).
    \item Các lớp tương tác cụ thể kế thừa từ \texttt{MonoBehaviour} và quản lý va chạm thông qua các trigger vật lý 2D. \texttt{ItemPickup} (quản lý nhặt vật phẩm và lưu trạng thái thu thập qua \texttt{WorldPickupSaveService}), \texttt{SpellBookChest} (mở hòm sách phép thuật), \texttt{Storage} (mở rương quản lý túi đồ), \texttt{ElevatorLever} (gạt cần kích hoạt thang máy), và \texttt{InteractPromptTrigger} (hiển thị gợi ý phím bấm).
    \item Các lớp tương tác này phụ thuộc trực tiếp vào \texttt{UIManager} để gửi thông điệp hiển thị phím gợi ý tương tác (\texttt{ShowInteractPrompt}) lên Canvas HUD của người chơi.
\end{itemize}

\begin{figure}[H]
    \centering
    \includegraphics[width=0.8\textwidth]{Hinhve/enemy_class_hierarchy.png}
    \caption{Biểu đồ phân cấp kế thừa các thực thể Kẻ địch (Enemy Class Hierarchy)}
    \label{fig:enemy_class_hierarchy}
\end{figure}

\textit{Giải thích biểu đồ kế thừa kẻ địch (Hình~\ref{fig:enemy_class_hierarchy}):}
Biểu đồ thể hiện mối quan hệ kế thừa trực tiếp từ lớp cơ sở trừu tượng \texttt{MvEnemyBase}. Các thực thể quái vật thông thường (\texttt{MvEm0010} - BoneDog, \texttt{MvEm0060} - ShieldKnight, \texttt{MvEm0070} - FlyingMage, \texttt{MvEm0100} - Brawler) cùng các Boss lớn có AI phức tạp hơn (\texttt{MvEm9020} - Golem Boss, \texttt{MvEm9030} - Marionette Boss) đều kế thừa trực tiếp từ lớp cơ sở này để đảm bảo tính đồng nhất của hệ thống vật lý và nhận sát thương.


\subsection{Cấu trúc chi tiết các phân hệ cốt lõi}
\label{subsection:5.2.4}

\subsubsection{Phân hệ Người chơi và Cơ chế Biến hình (Form Switch)}
Cơ chế biến hình cho phép người chơi thay đổi linh hoạt hình dạng của nhân vật chính sang các dạng quái vật khác khi thu thập linh hồn. Hệ thống này sử dụng cấu trúc dữ liệu tĩnh hướng cấu hình (Data-Driven Design) thông qua các lớp:
\begin{itemize}
    \item \texttt{PlayerFormId} (Enum): Định danh các hình thái linh hồn: \texttt{Human} (Dạng người mặc định), \texttt{Em0010} (BoneDog), \texttt{Em0020} (GhostCat), \texttt{Em0060} (ShieldKnight), \texttt{Em0070} (FlyingMage), và \texttt{Em0100} (Brawler).
    \item \texttt{PlayerFormDataSO} (ScriptableObject): Lưu trữ các dữ liệu thiết kế cố định như tên hiển thị, icon HUD, lượng HP tối đa, Animator Controller và Prefab vật lý tương ứng của từng hình thái.
\end{itemize}

Khi người chơi kích hoạt chuyển đổi linh hồn gần một xác chết kẻ địch, lớp điều phối trung tâm \texttt{PlayerController} kết hợp với \texttt{PlayerFormManager} thực hiện chuỗi hành vi đồng bộ hóa thời gian thực:
\begin{enumerate}
    \item \textbf{Đồng bộ hóa Chỉ số Sinh lực:} Lưu chỉ số HP hiện tại của form cũ và thiết lập chỉ số HP tương ứng của form mới vào lớp \texttt{PlayerStats}.
    \item \textbf{Tải cấu trúc vật lý động:} Lớp \texttt{PlayerFormManager} hủy phiên bản Prefab cũ và khởi tạo (\texttt{Instantiate}) Prefab của form mới làm con của Player Root. Prefab này được đính kèm component \texttt{PlayerFormBodyRef} định nghĩa rõ ràng các collider va chạm môi trường và các trigger hitbox tấn công cận chiến đặc thù.
    \item \textbf{Đồng bộ hóa Collider trong Physics Engine:} Lớp \texttt{PlayerMovement} tiếp nhận thông tin từ \texttt{PlayerFormBodyRef} mới để gán lại Collider va chạm thực tế và các điểm kiểm tra mặt đất (\texttt{GroundCheck}), vách tường (\texttt{WallCheck}) vào Physics2D của Unity, tránh lỗi kẹt địa hình do chênh lệch kích thước giữa các hình thái.
    \item \textbf{Rebind hoạt ảnh:} Thay đổi runtime controller của thành phần \texttt{Animator} trên Player Root bằng Animator Controller của form mới để kích hoạt tức thì các clip hoạt ảnh tương ứng. Hệ thống quản lý hoạt ảnh của nhân vật chính và các quái vật thông thường thực hiện gọi trực tiếp các hoạt ảnh này thông qua chuỗi ký tự (string-based loading) thay vị sử dụng các tham số và quan hệ chuyển cảnh phức tạp trong Animator Controller.
    \item \textbf{Cập nhật chỉ số chiến đấu:} Thay đổi cấu hình sát thương và chuỗi combo cận chiến trong \texttt{PlayerCombat} theo thông số của form mới.
\end{enumerate}

Để tránh việc lớp \texttt{PlayerController} bị phình to do phải quản lý hàng chục trạng thái của nhiều form khác nhau, hệ thống sử dụng mẫu thiết kế Factory thông qua lớp \texttt{FormStateFactory}. Khi người chơi ở dạng người mặc định, \texttt{PlayerStateMachine} sử dụng các trạng thái tĩnh được khởi tạo từ đầu (như \texttt{IdleState}, \texttt{MoveState}, \texttt{JumpState}, \texttt{DashState}, \texttt{WallClimbState}). Khi chuyển sang một dạng Enemy biến hình, \texttt{FormStateFactory} sẽ tạo ra theo nhu cầu (on-demand) một bộ trạng thái chuyên biệt cho form đó (ví dụ: \texttt{Em0010IdleState}, \texttt{Em0010AttackState}) kế thừa từ lớp cơ sở \texttt{PlayerState}.

\subsubsection{Phân hệ Kẻ địch và Trí tuệ nhân tạo (AI Brain)}
Hệ thống quái vật trong game được thiết kế hướng đối tượng chặt chẽ nhằm tối ưu hóa việc quản lý mã nguồn di chuyển, phát hiện mục tiêu và va chạm vật lý:
\begin{itemize}
    \item \textbf{Lớp cơ sở trừu tượng \texttt{MvEnemyBase}:} Kế thừa trực tiếp từ \texttt{MonoBehaviour} và thực thi giao diện \texttt{IDamageable}. Lớp này chịu trách nhiệm quản lý máy trạng thái \texttt{EnemyStateMachine}, quản lý các chỉ số cơ bản (Máu, Tốc độ, Tầm đánh), xử lý các hiệu ứng vật lý (lực đẩy lùi - Knockback, đệm va chạm), và tích hợp hệ thống tự động tìm kiếm mục tiêu (\texttt{UpdateSearchTarget}) trong phạm vi quạt góc quét (\texttt{searchFanAngle}) và tầm quét (\texttt{searchRadius}). Đồng thời, lớp này khai báo cấu trúc enum \texttt{AsCommon} đại diện cho các trạng thái dùng chung của quái vật (Idle, Run, Turn, Hit, Death).
    \item \textbf{Các lớp quái vật và Boss cụ thể (\texttt{MvEmXXXX}):} Kế thừa từ \texttt{MvEnemyBase} (ví dụ: \texttt{MvEm0010} - BoneDog, \texttt{MvEm9020} - Golem Boss, \texttt{MvEm9030} - Marionette Boss). Các lớp này override các phương thức đặc thù như thiết lập kiểu quái vật (\texttt{EmType}), ánh xạ tên trạng thái, đăng ký bảng trạng thái và tùy biến thuật toán quét tìm người chơi.
\end{itemize}

Các trận chiến với Boss lớn được điều phối bởi lớp bộ não AI chuyên biệt \texttt{MvEmBrain} chạy song song với máy trạng thái:
\begin{itemize}
    \item \textbf{Trọng số mong muốn (Desire Weights):} AI liên tục đánh giá khoảng cách đến người chơi và phần trăm máu còn lại của Boss để cập nhật bảng trọng số mong muốn của các chiêu thức (ví dụ: khi người chơi ở xa, trọng số chiêu phóng đạn tăng; khi Boss dưới 50\% máu, trọng số các đòn tấn công nhanh, diện rộng tăng).
    \item \textbf{Chuỗi chiêu thức (Combo Patterns):} Thay vì ra đòn ngẫu nhiên, Boss chọn một chuỗi hành động kết hợp được lập trình sẵn (\texttt{ComboBuff} lưu mảng các trạng thái hành vi). Khi combo bắt đầu, Boss sẽ thực thi tuần tự tất cả các đòn đánh trong combo đó trước khi quay lại trạng thái đánh giá mong muốn tiếp theo. Sơ đồ luồng quyết định hành động của Boss AI được mô tả chi tiết tại Hình~\ref{fig:boss_ai_flow}.
\end{itemize}

\begin{figure}[H]
    \centering
    \includegraphics[height=10cm]{Hinhve/BossAiFlow.png}
    \caption{Sơ đồ luồng ra quyết định hành động của Boss AI dựa trên Desire và Combo}
    \label{fig:boss_ai_flow}
\end{figure}

\textit{Giải thích Sơ đồ luồng quyết định của Boss AI (Hình~\ref{fig:boss_ai_flow}):}
Sơ đồ minh họa quá trình tính toán và ra quyết định hành vi của bộ não Boss AI:
\begin{enumerate}
    \item Hệ thống AI liên tục đo đạc khoảng cách vật lý từ vị trí Boss đến người chơi và phần trăm lượng máu còn lại của Boss.
    \item Từ các chỉ số runtime này, AI thực hiện cập nhật bảng trọng số mong muốn (Desire Weights) của các chiêu thức có sẵn trong cơ sở dữ liệu.
    \item Nếu Boss đang không thực thi một chuỗi chiêu thức cụ thể, AI sẽ lựa chọn chiêu thức có trọng số mong muốn cao nhất để thi triển.
    \item Nếu chiêu thức được chọn là một đòn đánh đơn lẻ, Boss sẽ thực thi đòn đánh và quay lại pha đánh giá. Nếu chiêu thức kích hoạt một chuỗi combo (Combo Pattern), Boss sẽ chạy tuần tự toàn bộ mảng trạng thái hành vi được lưu trong \texttt{ComboBuff} (ví dụ: lướt tới áp sát, vung kiếm quét rộng, rồi đập đất tạo chấn động) trước khi cho phép AI thực hiện pha đánh giá tiếp theo.
\end{enumerate}

\subsubsection{Các hệ thống lõi phục vụ vận hành}
\begin{itemize}
    \item \textbf{Hệ thống Lưu trữ (Save/Load):} Sử dụng lớp \texttt{GameSaveManager} để serialize trạng thái trò chơi thành chuỗi dữ liệu JSON và lưu trữ lâu dài. \texttt{GameSaveManager} được thiết kế theo mẫu Singleton và nằm trên Object có thuộc tính \texttt{DontDestroyOnLoad}, đảm bảo tiến trình lưu/tải không bị gián đoạn khi chuyển cảnh. Để làm rõ các tương tác truyền tin trong quá trình lưu trữ, sơ đồ trình tự quá trình lưu trữ được đặc tả tại Hình~\ref{fig:save_system_sequence}.
\end{itemize}

\begin{figure}[H]
    \centering
    \includegraphics[width=0.85\textwidth]{Hinhve/save_system_sequence.png}
    \caption{Sơ đồ trình tự quá trình lưu trữ tiến trình game (Save Game Sequence)}
    \label{fig:save_system_sequence}
\end{figure}

\textit{Giải thích Sơ đồ trình tự quá trình lưu trữ (Hình~\ref{fig:save_system_sequence}):}
Biểu đồ trình tự thể hiện sự tương tác truyền thông điệp giữa các đối tượng khi người chơi thực hiện lưu game:
\begin{enumerate}
    \item Quá trình lưu game có thể được kích hoạt qua hai luồng:
    \begin{itemize}
        \item \textbf{Lưu thủ công (Manual Save):} Người chơi di chuyển lại gần Đền thờ hồi sinh (\texttt{RespawnShrine}) và nhấn phím tương tác \texttt{E}. Lớp \texttt{RespawnShrine} nhận tín hiệu và gọi trực tiếp \texttt{SaveActiveSlot()} trên \texttt{GameSaveManager}.
        \item \textbf{Lưu tự động (Auto Save):} Khi người chơi thực hiện các sự kiện quan trọng (như mở khóa cổng, nâng cấp kỹ năng, mua vật phẩm), các hệ thống tương ứng (ví dụ: \texttt{PortalPoint}, \texttt{SkillTreeManager}) sẽ gửi yêu cầu \texttt{Request()} tới lớp tĩnh \texttt{GameAutoSave}. Lớp này sau đó sẽ gọi \texttt{SaveActiveSlot()} của \texttt{GameSaveManager}.
    \end{itemize}
    \item \texttt{GameSaveManager} kích hoạt hàm nội bộ \texttt{CaptureSaveData()} để tiến hành thu thập dữ liệu hiện thời từ các phân hệ.
    \item \texttt{GameSaveManager} gửi yêu cầu lấy số dư tiền vàng từ ví \texttt{CurrencyWalletSO}, lấy thông tin vật phẩm từ \texttt{InventoryStateSO} và cấp độ nâng cấp từ \texttt{FacilityManager}.
    \item Sau khi thu thập đầy đủ, \texttt{GameSaveManager} tiến hành serialize đối tượng \texttt{GameSaveData} thành chuỗi văn bản định dạng JSON.
    \item \texttt{GameSaveManager} gọi thư viện hệ thống để ghi đè hoặc tạo mới file JSON lưu trữ cục bộ xuống ổ cứng.
    \item Khi ổ cứng phản hồi ghi tệp tin thành công, đối với luồng lưu thủ công, \texttt{RespawnShrine} sẽ kích hoạt hiển thị thông báo lưu game thành công lên màn hình HUD.
\end{enumerate}

\begin{itemize}
    \item \textbf{Hệ thống Culling tối ưu hóa hiệu năng (Culling System):} Trò chơi áp dụng cơ chế Culling linh hoạt thông qua sự kết hợp giữa culling theo phòng chơi (\texttt{RoomController}) và culling theo khoảng cách (\texttt{DistanceCullingTarget}), được điều phối bởi \texttt{CullingManager}. Thay vì sử dụng \texttt{SetActive(false)} thô bạo lên toàn bộ GameObject có thể gây ra chi phí phân bổ bộ nhớ mỗi khi bật/tắt, hệ thống tiến hành vô hiệu hóa chọn lọc các component nặng tài nguyên xử lý. Cụ thể, các component như script AI (\texttt{MvEnemyBase}), script xử lý hành vi (\texttt{MvAttack}), khóa cứng vật lý của \texttt{Rigidbody2D}, dừng hoạt ảnh \texttt{Animator} và tạm dừng hệ thống hạt \texttt{ParticleSystem} sẽ bị vô hiệu hóa khi đối tượng thỏa mãn điều kiện culling, giúp giảm thiểu tải xử lý CPU mà không làm xáo trộn cấu trúc GameObject (Hình~\ref{fig:culling_vis}).
\end{itemize}

\begin{figure}[H]
    \centering
    \includegraphics[width=0.8\textwidth]{Hinhve/CullingVisualization.png}
    \caption{Hình vẽ trực quan các phòng chơi bị tắt hiển thị ngoài tầm nhìn để tối ưu hóa}
    \label{fig:culling_vis}
\end{figure}

\textit{Giải thích cơ chế tối ưu hóa culling phòng chơi và khoảng cách (Hình~\ref{fig:culling_vis}):}
Hình vẽ trực quan minh họa giải thuật tối ưu hóa hiệu năng tính toán vật lý và logic thông qua culling các vùng không hoạt động:
\begin{enumerate}
    \item Bản đồ thế giới game được chia nhỏ thành các khu vực phòng độc lập, mỗi phòng được quản lý bởi một \texttt{RoomController} định nghĩa bằng Trigger Collider. Các thực thể động bên trong phòng sẽ đóng vai trò là một \texttt{RoomCullingMember}.
    \item Khi người chơi di chuyển bên trong phòng, phòng đó được đánh dấu là thức dậy. Khi người chơi rời khỏi, phòng bị đánh dấu là ngủ đông (Sleeping). Ở trạng thái ngủ đông, \texttt{RoomCullingMember} sẽ ép đối tượng vào trạng thái vô hiệu hóa các component và khóa không cho \texttt{DistanceCullingTarget} tự ý kích hoạt lại đối tượng, đảm bảo quái vật không thức giấc ngoài ý muốn.
    \item Khi ở trong một phòng đang hoạt động, \texttt{DistanceCullingTarget} tiếp tục đảm nhiệm chức năng vi chỉnh: tính toán khoảng cách từ người chơi đến đối tượng. Nếu đối tượng nằm ngoài bán kính vô hiệu hóa mặc định, nó vẫn sẽ bị tắt các component nhằm tối ưu tối đa lượng tính toán của AI quái vật hay đạn bay ở xa, giúp duy trì hiệu năng 60 FPS ổn định trên cấu hình tối thiểu.
\end{enumerate}

\begin{itemize}
    \item \textbf{Hệ thống Nâng cấp và Cửa hàng:} \texttt{FacilityManager} chịu trách nhiệm đọc ghi tiến trình nâng cấp chỉ số vĩnh viễn (HP, Mana, Sát thương cận chiến) dựa trên số lượng các vật phẩm đặc biệt. \texttt{ShopStateSO} sử dụng ScriptableObject để lưu trữ danh sách các kỹ năng phép thuật và công cụ bổ trợ có bán tại NPC Shopkeeper, tự động đồng bộ hóa trạng thái đã mua/chưa mua giữa các màn chơi.
\end{itemize}

\subsubsection{Giao diện người dùng và Canvas (UI Canvas Architecture)}
Hệ thống giao diện người dùng (UI) được thiết kế dựa trên cấu trúc phân lớp Canvas của Unity Engine:
\begin{itemize}
    \item \textbf{HUD Canvas (Screen Space - Overlay):} Quản lý các thông tin động thời gian thực hiển thị liên tục (HP, Mana, Stamina, ô phép thuật, tiền vàng). Canvas này sử dụng thành phần \texttt{Canvas Scaler} chế độ \texttt{Scale With Screen Size} với độ phân giải tham chiếu là \texttt{1920x1080}.
    \item \textbf{Menu \& Panels Canvas (Screen Space - Overlay):} Quản lý các giao diện tĩnh (Inventory, Status, Shop, Facility, Full Map). Giao diện này được điều khiển bởi \texttt{UIManager} và \texttt{UIStateManager} để kiểm soát vòng đời và độ ưu tiên hiển thị.
    \item \textbf{World Space Canvas:} Sử dụng riêng biệt cho các thành phần UI động cần kết nối trực tiếp với tọa độ thế giới 2D, tiêu biểu là lớp hiển thị chữ sát thương bay \texttt{DamageTextView} xuất hiện ngay trên đầu các thực thể khi nhận sát thương chiến đấu.
\end{itemize}

Để thực hiện cơ chế tạm dừng trò chơi khi mở các bảng chức năng tĩnh, hệ thống sử dụng lớp \texttt{UIPauseRequester} gửi yêu cầu tới bộ quản lý trung tâm \texttt{GamePauseManager} để thiết lập \texttt{Time.timeScale = 0f}. Việc này dừng hoàn toàn mọi chuyển động vật lý của Rigidbody2D, các script tính toán AI quái vật và hệ thống đạn phép, đảm bảo an toàn cho nhân vật chính trong khi người chơi tương tác với giao diện.


\section{Thiết kế chi tiết}
\label{section:5.3}

\subsection{Thiết kế lớp}
\label{subsection:5.3.1}
Để minh họa cụ thể cấu trúc lập trình và thiết kế chi tiết thuộc tính, phương thức của 3 lớp core xương sống đóng vai trò cốt lõi trong hệ thống (\texttt{PlayerController}, \texttt{MvEnemyBase}, và \texttt{GameSaveManager}).

Dưới đây là biểu đồ lớp chi tiết của ba lớp được kết xuất trực tiếp từ mã nguồn thực tế:

\begin{figure}[H]
    \centering
    \includegraphics[width=0.45\textwidth]{Hinhve/player_controller_class.png}
    \caption{Biểu đồ lớp chi tiết của lớp PlayerController}
    \label{fig:player_controller_class}
\end{figure}

\textit{Giải thích biểu đồ lớp PlayerController (Hình~\ref{fig:player_controller_class}):}
Lớp \texttt{PlayerController} quản lý và phối hợp hoạt động của thực thể nhân vật chính. Các thuộc tính và phương thức được đặc tả như sau:
\begin{itemize}
    \item \textbf{Các thuộc tính thành viên:}
    \begin{itemize}
        \item \texttt{stateMachine}: Đối tượng kiểu \texttt{PlayerStateMachine} điều khiển vòng đời và việc chuyển đổi trạng thái di chuyển hiện thời của nhân vật.
        \item \texttt{playerInput}: Tham chiếu đến component xử lý nhập liệu từ bàn phím và chuột của người dùng.
        \item \texttt{playerMovement}: Thành phần tính toán động lực học vật lý, chịu trách nhiệm áp dụng lực lên Rigidbody2D.
        \item \texttt{playerStats}: Đối tượng lưu trữ các chỉ số sinh mệnh runtime của nhân vật (Máu, Thể lực, Năng lượng).
        \item \texttt{playerFormManager}: Trình quản lý cấu trúc prefab vật lý của hình thái linh hồn khi chuyển dạng.
        \item \texttt{currentForm}: Enum xác định nhân vật đang ở dạng người hay dạng biến hình.
    \end{itemize}
    \item \textbf{Các phương thức xử lý:}
    \begin{itemize}
        \item \texttt{TakeDamage(float damage)}: Phương thức thực thi từ giao diện \texttt{IDamageable} nhận lượng sát thương truyền vào, trừ trực tiếp máu trong \texttt{playerStats} và kích hoạt hiệu ứng rung màn hình.
        \item \texttt{TriggerTrapRespawn()}: Xử lý khi nhân vật chạm bẫy nhưng vẫn còn sống. Dịch chuyển nhân vật về vị trí an toàn trên mặt đất gần nhất, chỉ trừ sát thương bẫy mà không tải lại màn chơi.
        \item \texttt{PrepareRespawnFromDeathSequenceRoutine()}: Thực thi chuỗi hồi sinh khi người chơi mất toàn bộ sinh lực, tự động phục hồi hoàn toàn các chỉ số, tải lại màn chơi và ép buộc nhân vật quay về điểm lưu trữ mặc định (Nhà thờ - Church).
        \item \texttt{AbandonToShrine()}: Kích hoạt sự kiện đầu hàng, vô hiệu hóa điều khiển và tải lại màn chơi để đưa người chơi quay về khu vực xuất phát (Church), tương tự như khi tử vong.
        \item \texttt{ToggleTransformForm()}: Phương thức kiểm tra điều kiện an toàn và kích hoạt việc biến đổi hình dạng linh hồn tại xác quái vật hoặc hủy chuyển dạng để quay về dạng người.
        \item \texttt{SyncPlayerMovementCollider()}: Phương thức tính toán lại kích thước và vị trí của collider va chạm động trên \texttt{PlayerMovement} để khớp chính xác với hình thể của dạng quái vật vừa biến hình.
    \end{itemize}
\end{itemize}

\begin{figure}[H]
    \centering
    \includegraphics[width=0.35\textwidth]{Hinhve/enemy_base_class.png}
    \caption{Biểu đồ lớp chi tiết của lớp MvEnemyBase}
    \label{fig:enemy_base_class}
\end{figure}

\textit{Giải thích biểu đồ lớp MvEnemyBase (Hình~\ref{fig:enemy_base_class}):}
Lớp \texttt{MvEnemyBase} là lớp cơ sở trừu tượng định nghĩa khung hành vi và cấu trúc AI cho toàn bộ quái vật và Boss. Các thuộc tính và phương thức được đặc tả như sau:
\begin{itemize}
    \item \textbf{Các thuộc tính thành viên:}
    \begin{itemize}
        \item \texttt{maxHealth}: Lượng máu tối đa của quái vật được cấu hình tĩnh.
        \item \texttt{moveSpeed}: Tốc độ tuần tra cơ sở của thực thể.
        \item \texttt{enemyStateMachine}: Đối tượng máy trạng thái quản lý các trạng thái tuần tra, truy đuổi của quái vật.
        \item \texttt{target}: Đối tượng \texttt{Transform} tham chiếu trực tiếp đến nhân vật chính đang bị truy đuổi.
        \item \texttt{searchRadius}: Bán kính quét xung quanh thực thể để phát hiện người chơi.
        \item \texttt{searchFanAngle}: Góc quét hình quạt phía trước mặt quái vật để nhận diện mục tiêu.
        \item \texttt{isDead}: Biến logic xác định trạng thái sống/chết của quái vật.
    \end{itemize}
    \item \textbf{Các phương thức xử lý:}
    \begin{itemize}
        \item \texttt{TakeDamage(float damage)}: Phương thức xử lý khi nhận sát thương từ người chơi, trừ máu quái vật, kích hoạt hiệu ứng chữ sát thương bay (Damage Text) và đẩy lùi vật lý.
        \item \texttt{Die()}: Phương thức kích hoạt khi máu về 0, thực hiện tắt toàn bộ va chạm vật lý, chạy hiệu ứng tan biến và gọi sự kiện giải phóng bộ nhớ.
        \item \texttt{UpdateSearchTarget()}: Chạy định kỳ để thực hiện quét va chạm tròn (OverlapCircle), xác định xem nhân vật chính có nằm trong tầm nhìn và góc quét hay không để cập nhật thuộc tính \texttt{target}.
        \item \texttt{MovePatrol()}: Phương thức thực hiện di chuyển tuần tra qua lại quanh vị trí xuất phát ban đầu dựa trên một khoảng cách bán kính định sẵn.
        \item \texttt{MoveReturnToOrigin()}: Ra lệnh cho quái vật quay trở lại vị trí đền thờ hoặc điểm tuần tra ban đầu khi người chơi di chuyển ra khỏi tầm truy đuổi.
        \item \texttt{ApplyKnockback()}: Thực hiện tác dụng lực đẩy lùi vật lý lên Rigidbody2D của quái vật theo hướng ngược lại với nguồn gây sát thương.
    \end{itemize}
\end{itemize}

\begin{figure}[H]
    \centering
    \includegraphics[width=0.35\textwidth]{Hinhve/save_manager_class.png}
    \caption{Biểu đồ lớp chi tiết của lớp GameSaveManager}
    \label{fig:save_manager_class}
\end{figure}

\textit{Giải thích biểu đồ lớp GameSaveManager (Hình~\ref{fig:save_manager_class}):}
Lớp \texttt{GameSaveManager} chịu trách nhiệm lưu trữ và phục hồi tiến trình chơi của game. Các thuộc tính và phương thức được đặc tả như sau:
\begin{itemize}
    \item \textbf{Các thuộc tính thành viên tiêu biểu:}
    \begin{itemize}
        \item \texttt{instance}: Tham chiếu tĩnh duy nhất trỏ tới đối tượng quản lý đang hoạt động trong bộ nhớ (Singleton Pattern).
        \item \texttt{saveFileName}: Chuỗi ký tự định nghĩa tên tệp lưu trữ mặc định trên thiết bị.
        \item \texttt{activeSlotIndex}: Chỉ số ô lưu trữ đang được nạp để chơi game.
        \item \textbf{Các cơ sở dữ liệu (Databases):} \texttt{itemDatabase}, \texttt{spellDatabase} chứa cấu hình tài nguyên tĩnh của game.
        \item \textbf{Các trạng thái chạy (Runtime States):} \texttt{currencyWallet} (tiền vàng), \texttt{inventoryState} (túi đồ), \texttt{toolEquip} (công cụ), \texttt{healingPotionEquip} (bình máu), \texttt{shopState} (cửa hàng), \texttt{skillProgress} (tiến trình kỹ năng), \texttt{spellEquip} (phép thuật), \texttt{facilityManager} (nâng cấp làng), \texttt{skillTreeManager} (nhánh kỹ năng). Các tham chiếu này hỗ trợ \texttt{GameSaveManager} liên kết để lấy hoặc nạp dữ liệu từ các phân hệ.
    \end{itemize}
    \item \textbf{Các phương thức xử lý:}
    \begin{itemize}
        \item \texttt{SaveGame()} / \texttt{SaveActiveSlot()} / \texttt{SaveGameToSlot()}: Các phương thức chịu trách nhiệm chuyển đổi và lưu toàn bộ tiến trình của ô nhớ hiện tại hoặc được chỉ định xuống tệp tin.
        \item \texttt{LoadGame()} / \texttt{LoadGameFromSlot()}: Đọc tệp tin JSON tương ứng với ô nhớ được chọn, giải mã ngược sang cấu trúc dữ liệu đối tượng trong C\#.
        \item \texttt{CreateNewGameSlot()}: Tạo ra một ô lưu trữ trống mới với các giá trị thời gian và chỉ số khởi tạo ban đầu (New Game).
        \item \texttt{CaptureSaveData()}: Phương thức thu thập các chỉ số thực tế từ tất cả các ScriptableObjects và Manager hiện hành, gom nhóm lại vào đối tượng cấu trúc trung gian \texttt{GameSaveData} để chuẩn bị serialize JSON.
        \item \texttt{ApplySaveData()}: Lấy đối tượng \texttt{GameSaveData} đã được phân giải, sau đó phân phối lại các luồng dữ liệu tương ứng nạp vào các hệ thống \texttt{Runtime States} khi tải game.
        \item \texttt{DeleteSave()}: Xóa tệp tin lưu trữ ra khỏi ổ cứng.
    \end{itemize}
\end{itemize}

Nhằm minh họa luồng truyền thông điệp giữa các đối tượng trong hệ thống, sơ đồ tuần tự (Hình~\ref{fig:form_switch_sequence}) đặc tả chi tiết quy trình tương tác của use case Chuyển đổi linh hồn (Form Switch) giữa các thành phần: \texttt{PlayerInput}, \texttt{PlayerController}, \texttt{PlayerFormManager}, và \texttt{PlayerStats}.
\begin{figure}[H]
    \centering
    \includegraphics[width=0.85\textwidth]{Hinhve/FormSwitchSequence.png}
    \caption{Sơ đồ trình tự quá trình đồng bộ hóa chuyển đổi hình thái (Form Switch Sequence)}
    \label{fig:form_switch_sequence}
\end{figure}

\textit{Giải thích Sơ đồ trình tự chuyển đổi hình thái linh hồn (Hình~\ref{fig:form_switch_sequence}):}
Biểu đồ mô tả chi tiết chuỗi tương tác thời gian thực giữa các thành phần con của thực thể người chơi khi thực hiện cơ chế biến hình (Form Switch):
\begin{enumerate}
    \item Khi ở gần xác các Enemy đặc biệt, người chơi nhấn phím tương tác biến hình (\texttt{Tab} hoặc phím gán). \texttt{PlayerInput} nhận diện tín hiệu và gửi thông báo sự kiện nhập liệu đến bộ điều phối trung tâm \texttt{PlayerController}.
    \item \texttt{PlayerController} đóng vai trò là Mediator kiểm tra điều kiện an toàn (không bị khống chế, không đang rơi tự do) và gọi phương thức \texttt{SwitchForm()} đến \texttt{PlayerFormManager}.
    \item \texttt{PlayerFormManager} thực hiện truy xuất thông tin cấu hình hình thái linh hồn mới từ tài sản tĩnh ScriptableObject \texttt{PlayerFormDataSO}.
    \item \texttt{PlayerFormManager} tiến hành đồng bộ hóa các chỉ số máu tối đa, thể lực tối đa mới vào component quản lý thuộc tính \texttt{PlayerStats}.
    \item \texttt{PlayerFormManager} thực hiện hủy đối tượng Prefab hình thể cũ và khởi tạo động (\texttt{Instantiate}) Prefab của hình thái linh hồn mới làm con của Player Root.
    \item Lớp điều phối \texttt{PlayerController} gọi hàm \texttt{SyncPlayerMovementCollider()} đến component \texttt{PlayerMovement} để gán lại kích thước Collider vật lý và các điểm GroundCheck/WallCheck khớp chính xác với hình dạng mới, ngăn ngừa hiện tượng lỗi kẹt địa hình. Đồng thời, Animator của hình thái mới được rebind vào hệ thống hoạt ảnh để chạy mượt mà chuỗi chuyển động tương ứng.
\end{enumerate}


\subsection{Thiết kế cơ sở dữ liệu}
\label{subsection:5.3.2}
Vì trò chơi được thiết kế như một phần mềm độc lập chạy ngoại tuyến hoàn toàn trên PC, hệ thống không sử dụng các hệ quản trị cơ sở dữ liệu quan hệ máy chủ (như SQL Server hay MySQL) để tránh đòi hỏi kết nối mạng và các chi phí cài đặt phức tạp cho người chơi.

Thay vào đó, trò chơi áp dụng thiết kế cơ sở dữ liệu hướng tài liệu cục bộ (document-based local file database) bằng định dạng JSON bền vững. Cơ sở dữ liệu được tổ chức dưới dạng một lớp tổng hợp gốc là \texttt{GameSaveData} chứa các tập hợp dữ liệu con đại diện cho trạng thái của từng phân hệ. Việc tổ chức cơ sở dữ liệu dưới dạng JSON giúp đảm bảo tốc độ đọc ghi cực nhanh, dung lượng lưu trữ nhẹ (dưới 100KB mỗi slot lưu trữ) và dễ dàng sao lưu, phục hồi thông qua các thư viện hệ thống tệp chuẩn của hệ điều hành.

Hệ thống quản lý lưu trữ được phân tách thành ba dạng độc lập: lưu trữ cấu hình âm thanh hệ thống bằng cơ chế \texttt{PlayerPrefs} của Unity, lưu trữ thiết lập phím bấm (Keybinds) người dùng bằng định dạng JSON độc lập, và lưu trữ tiến trình chơi game bằng định dạng JSON bền vững. Các thuộc tính và thành phần dữ liệu được quy định chi tiết bằng cấu trúc danh sách đặc tả dưới đây:

\textbf{Lưu trữ cấu hình hệ thống (qua PlayerPrefs):}
Các thiết lập âm lượng hệ thống cơ bản được ghi nhận cục bộ dưới dạng cặp khóa - giá trị bao gồm:
\begin{itemize}
    \item \texttt{Pref\_MasterVolume} (kiểu dữ liệu \texttt{float}, giá trị mặc định \texttt{1.0f}): Âm lượng tổng của toàn bộ trò chơi (Master Volume).
    \item \texttt{Pref\_BgmVolume} (kiểu dữ liệu \texttt{float}, giá trị mặc định \texttt{0.8f}): Âm lượng của phần âm nhạc nền (BGM Volume).
    \item \texttt{Pref\_SfxVolume} (kiểu dữ liệu \texttt{float}, giá trị mặc định \texttt{0.8f}): Âm lượng của phần hiệu ứng âm thanh chiến đấu (SFX Volume).
\end{itemize}

\textbf{Lưu trữ thiết lập phím bấm (Keybinds - qua input\_profile.json):}
Các phím điều khiển do người dùng tùy biến được lưu tại đường dẫn thư mục tài liệu cá nhân (\texttt{MyDocuments/DreamKnight/input\_profile.json}) dưới cấu trúc lớp \texttt{InputBindingProfile} chứa danh sách các phần tử \texttt{InputBindingEntry} bao gồm:
\begin{itemize}
    \item \texttt{action} (kiểu dữ liệu \texttt{string}): Tên hành động có thể gán phím, ví dụ: \texttt{MoveUp}, \texttt{MoveDown}, \texttt{MoveLeft}, \texttt{MoveRight}, \texttt{Jump}, \texttt{Dodge}, \texttt{NormalAttack}, \texttt{Interact}, \texttt{UsePotion}, \texttt{UseTool}, \texttt{UseSpell}, \texttt{Transform}.
    \item \texttt{controlPath} (kiểu dữ liệu \texttt{string}): Đường dẫn phím bấm của New Input System tương ứng được gán cho hành động đó, ví dụ: \texttt{<Keyboard>/w}, \texttt{<Keyboard>/space}, \texttt{<Mouse>/leftButton}.
\end{itemize}

\textbf{Lưu trữ tiến trình chơi game chi tiết (qua GameSaveData):}
Dữ liệu chi tiết về tiến trình chơi của người chơi được lưu trong lớp cấu trúc \texttt{GameSaveData} bao gồm các trường:
\begin{itemize}
    \item \texttt{gold} (kiểu dữ liệu \texttt{int}): Số lượng tiền vàng tích lũy hiện có của nhân vật chính.
    \item \texttt{inventory} (kiểu dữ liệu \texttt{InventorySaveData}): Danh sách lưu trữ thông tin các vật phẩm thu thập được trong túi đồ.
    \item \texttt{toolEquip} (kiểu dữ liệu \texttt{EquipmentSaveData}): Dữ liệu về trang bị công cụ (Tool) hiện tại của người chơi.
    \item \texttt{healingPotionEquip} (kiểu dữ liệu \texttt{EquipmentSaveData}): Dữ liệu về trang bị bình thuốc hồi máu hiện tại của người chơi.
    \item \texttt{shop} (kiểu dữ liệu \texttt{ShopSaveData}): Trạng thái đã mua hoặc chưa mua các vật phẩm đặc trưng trong cửa hàng NPC.
    \item \texttt{skillProgress} (kiểu dữ liệu \texttt{SkillProgressSaveData}): Dữ liệu theo dõi trạng thái mở khóa các kỹ năng phép thuật đã học.
    \item \texttt{facilityProgress} (kiểu dữ liệu \texttt{FacilityProgressSaveData}): Tiến trình nâng cấp vĩnh viễn các cơ sở thuộc tính cơ bản của nhân vật tại làng.
    \item \texttt{skillTreeProgress} (kiểu dữ liệu \texttt{SkillTreeProgressSaveData}): Tiến trình mở khóa các nhánh kỹ năng chiến đấu trong cây kỹ năng (Skill Tree).
    \item \texttt{portals} (kiểu dữ liệu \texttt{PortalSaveData}): Danh sách trạng thái kích hoạt của các cổng dịch chuyển nhanh (Portals).
    \item \texttt{doors} (kiểu dữ liệu \texttt{DoorSaveData}): Trạng thái đóng/mở của các cánh cửa lớn phân vùng bản đồ trong game.
    \item \texttt{worldPickups} (kiểu dữ liệu \texttt{WorldPickupSaveData}): Trạng thái ghi nhận các vật phẩm đã được nhặt trong màn chơi.
    \item \texttt{bossDefeats} (kiểu dữ liệu \texttt{BossDefeatSaveData}): Trạng thái đánh bại các Boss chính (đã tiêu diệt/chưa tiêu diệt) phục vụ cốt truyện.
    \item \texttt{equippedSpellId} (kiểu dữ liệu \texttt{string}): Định danh ID của kỹ năng phép thuật đang được người chơi chọn để thi triển.
    \item \texttt{slotIndex} (kiểu dữ liệu \texttt{int}): Chỉ mục ô nhớ lưu trữ game (Save Slot Index) đang được thực hiện.
    \item \texttt{playTimeSeconds} (kiểu dữ liệu \texttt{float}): Tổng thời gian chơi game đã ghi nhận thực tế (tính bằng giây).
\end{itemize}

\textbf{Phân loại các đối tượng ScriptableObjects trong kiến trúc game:}
Để hỗ trợ kiến trúc dữ liệu hướng cấu hình (Data-Driven Design), hệ thống phân tách ScriptableObjects thành hai nhóm riêng biệt:
\begin{itemize}
    \item \textbf{Nhóm dữ liệu cấu hình tĩnh (Static Configuration Assets):}
    \begin{itemize}
        \item \texttt{ItemDatabaseSO} (kế thừa từ \texttt{ScriptableObject}): Chứa danh sách định nghĩa thông số tĩnh của toàn bộ các vật phẩm trong trò chơi.
        \item \texttt{SpellDatabaseSO} (kế thừa từ \texttt{ScriptableObject}): Chứa danh sách định nghĩa thông số cơ bản của các kỹ năng phép thuật có thể học.
        \item \texttt{PlayerFormDataSO} (kế thừa từ \texttt{ScriptableObject}): Lưu trữ cấu hình prefab, icon HUD, lượng máu tối đa mặc định và Animator Controller của các hình thái linh hồn.
    \end{itemize}
    \item \textbf{Nhóm lưu trữ trạng thái chạy runtime (Runtime State Assets):}
    \begin{itemize}
        \item \texttt{InventoryStateSO} (kế thừa từ \texttt{ScriptableObject}): Quản lý danh sách các vật phẩm thực tế đang tồn tại trong túi đồ của người chơi.
        \item \texttt{CurrencyWalletSO} (kế thừa từ \texttt{ScriptableObject}): Quản lý số dư tiền vàng hiện thời và thực hiện các giao dịch cộng/trừ tiền trực tiếp.
        \item \texttt{ShopStateSO} (kế thừa từ \texttt{ScriptableObject}): Theo dõi và cập nhật trạng thái mua sắm vật phẩm và mở khóa kỹ năng của cửa hàng.
    \end{itemize}
\end{itemize}


\section{Xây dựng ứng dụng}
\label{section:5.4}

\subsection{Thư viện và công cụ sử dụng}
\label{subsection:5.4.1}
Quá trình xây dựng và phát triển trò chơi "Dream Witch" sử dụng các công cụ lập trình, thư viện bên thứ ba và phần mềm đồ họa với thông tin chi tiết được thống kê trong Bảng~\ref{table:development_tools}.

\begin{table}[H]
\centering
\begin{tabular}{|p{3.0cm}|p{4.5cm}|p{5.0cm}|}
\hline
\textbf{Mục đích sử dụng} & \textbf{Công cụ / Thư viện (Phiên bản)} & \textbf{Địa chỉ URL chính thức} \\ \hline
Engine phát triển game & Unity Editor (6000.0.40f1 LTS) & \url{https://unity.com/} \\ \hline
Trình soạn thảo mã nguồn & Visual Studio Code (v1.90+) & \url{https://code.visualstudio.com/} \\ \hline
Trợ lý lập trình AI & Antigravity (DeepMind Agent) & N/A (Tích hợp trong môi trường phát triển) \\ \hline
Hệ thống hội thoại & Yarn Spinner for Unity (v2.3.0) & \url{https://yarnspinner.dev/} \\ \hline
Diễn họa Boss 2D & Spine 2D Runtime (v4.2) & \url{https://esotericsoftware.com/} \\ \hline
Giải thuật tìm đường & A* Pathfinding Project (v4.2.17) & \url{https://arongranberg.com/astar/} \\ \hline
Quản lý phiên bản & Git \& GitHub Desktop & \url{https://github.com/} \\ \hline
Thiết kế đồ họa và Sprite & Adobe Photoshop \& Aseprite & \url{https://www.aseprite.org/} \\ \hline
\end{tabular}
\caption{Danh sách các thư viện và công cụ phát triển trò chơi}
\label{table:development_tools}
\end{table}

\subsection{Kết quả đạt được}
\label{subsection:5.4.2}
Sản phẩm sau quá trình phát triển là bản đóng gói trò chơi hoàn chỉnh hoạt động độc lập trên hệ điều hành Windows 64-bit (PC Standalone). Bản đóng gói bao gồm tệp thực thi \texttt{DreamWitch.exe} cùng thư mục tài nguyên \texttt{DreamWitch\_Data}. Mã nguồn C\# của dự án được tổ chức chặt chẽ theo cấu trúc phân rã thành các gói chức năng (packages) bao gồm: \texttt{Player}, \texttt{Enemy}, \texttt{System}, \texttt{UI}, \texttt{Interfaces}, \texttt{Core}, \texttt{Trap}, và \texttt{VFX}.

Để minh chứng cho quy mô phát triển của dự án, các thông tin định lượng chi tiết về mã nguồn C\# và dung lượng các thành phần liên quan được tổng hợp trong Bảng~\ref{table:project_stats}.

\begin{table}[H]
\centering
\begin{tabular}{|p{6.5cm}|p{2cm}|p{4cm}|}
\hline
\textbf{Chỉ số định lượng} & \textbf{Giá trị ghi nhận} & \textbf{Đơn vị đo lường} \\ \hline
Số lượng tệp tin mã nguồn C\# (\texttt{.cs}) & 347 & Tệp tin \\ \hline
Tổng số dòng mã nguồn (LOC) & 40.692 & Dòng code \\ \hline
Số lượng gói phân hệ chức năng chính & 8 & Gói (Namespace) \\ \hline
Dung lượng thư mục mã nguồn C\# thô & 1,61 & Megabytes (MB) \\ \hline
Dung lượng thư mục tài nguyên dự án (Assets) & 3,02 & Gigabytes (GB) \\ \hline
Dung lượng bản đóng gói sản phẩm (\texttt{.zip}) & $\sim$ 178 & Megabytes (MB) \\ \hline
\end{tabular}
\caption{Thống kê định lượng mã nguồn và tài nguyên dự án}
\label{table:project_stats}
\end{table}


\section{Sản phẩm}
\label{section:5.5}
Để đánh giá trực quan kết quả lập trình và thiết kế cơ chế, diễn tiến vận hành thực tế của trò chơi "Dream Witch" cần theo một kịch bản chơi game (gameplay scenario) hoàn chỉnh đi qua các phân khu bản đồ từ lúc khởi động cho tới khi hoàn thành nhiệm vụ cốt lõi.

Kịch bản trải nghiệm gồm 5 giai đoạn chính như sau:

\subsection{Giai đoạn 1: Khởi động, thiết lập ban đầu, khu vực xuất phát và khu vực nâng cấp Village}
Khi người chơi kích hoạt tệp \texttt{DreamWitch.exe}, hệ thống hiển thị màn hình khởi động (Hình~\ref{fig:titlescreen}). Giao diện cung cấp các tùy chọn để bắt đầu chơi mới hoặc tiếp tục tiến trình cũ thông qua giao diện chọn Save Slot (Hình~\ref{fig:mainmenu}). 

\begin{figure}[H]
    \centering
    \includegraphics[width=0.8\linewidth]{Hinhve/Title.jpg}
    \caption{Màn hình khởi động trò chơi (Title Screen)}
    \label{fig:titlescreen}
\end{figure}

\begin{figure}[H]
    \centering
    \includegraphics[width=0.8\linewidth]{Hinhve/SaveSlot.jpg}
    \caption{Giao diện chọn Slot (Save Slot Menu)}
    \label{fig:mainmenu}
\end{figure}
Sau khi vào game, nhân vật xuất hiện mặc định ở dạng Người (Human) tại khu vực xuất phát an toàn là Nhà thờ (Church) (Hình~\ref{fig:church}) để chuẩn bị cho hành trình thám hiểm. Người chơi có thể di chuyển tự do khám phá và tương tác với các NPC để nói chuyện hội thoại.
\begin{figure}[H]
    \centering
    \includegraphics[width=0.8\linewidth]{Hinhve/Church.jpg}
    \caption{Nhân vật tại khu vực xuất phát Nhà thờ (Church)}
    \label{fig:church}
\end{figure}
Tại khu vực Church, người chơi có thể tiến tới khu vực Village để nâng cấp trang bị và kĩ năng cho nhân vật cùng với mua những vật dụng cần thiết cho quá trình thám hiểm (Hình~\ref{fig:hud_start}).
\begin{figure}[H]
    \centering
    \includegraphics[width=0.8\linewidth]{Hinhve/Village.jpg}
    \caption{Khu vực làng (Village) nâng cấp và mua vật dụng}
    \label{fig:hud_start}
\end{figure}

\subsection{Giai đoạn 2: Thám hiểm khu vực Dungeon}
Nhân vật chính di chuyển qua các phân khu ngục tối (Dungeon) khác nhau là hai khu vực Ruin và Volcano. Người chơi thực hiện kiểm soát chuyển động nhân vật: chạy di chuyển mượt mà, thực hiện nhảy đơn, nhảy đôi để vượt qua các chướng ngại vật và bám vách leo tường vượt chướng ngại vật, cần phải vượt qua khu vực Ruin đầu tiên (Hình~\ref{fig:ruin_level}) và đánh bại Boss Ruin, trước khi mở khóa tiến sâu vào phân khu nguy hiểm hơn là khu vực Núi lửa (Volcano) (Hình~\ref{fig:volcano_level}).

\begin{figure}[H]
    \centering
    \includegraphics[width=0.8\linewidth]{Hinhve/RuinMap.jpg}
    \caption{Thám hiểm khu vực Di tích (Ruin)}
    \label{fig:ruin_level}
\end{figure}

\begin{figure}[H]
    \centering
    \includegraphics[width=0.8\linewidth]{Hinhve/VolcanoMap.jpg}
    \caption{Thám hiểm khu vực Núi lửa (Volcano)}
    \label{fig:volcano_level}
\end{figure}

\subsection{Giai đoạn 3: Cơ chế chiến đấu cận chiến và né tránh}
Trong quá trình thám hiểm , nhân vật đối đầu với quái vật tuần tra trên đường. Người chơi thi triển chuỗi đòn đánh cận chiến liên hoàn 3-hit combo tạo ra sát thương. Khi kẻ địch lao tới tấn công, người chơi kích hoạt phím lướt (Dash) để né tránh đòn đánh ngay trước thời điểm bị va chạm để không nhận sát thương (Hình~\ref{fig:quaivatthuong}).

\begin{figure}[H]
    \centering
    \includegraphics[width=0.8\linewidth]{Hinhve/Combat.jpg}
    \caption{Nhân vật cận chiến chém quái thú thường kèm chữ hiển thị sát thương bay}
    \label{fig:quaivatthuong}
\end{figure}

\subsection{Giai đoạn 4: Biến dạng thực thể và phép thuật tầm xa}
Trong quá trình thám hiểm sẽ có những quái vật đặc biệt khi đánh bại thì xác sẽ không biến mất, nếu người chơi đã mở node 1 của cây Skill Tree sẽ có thể thực hiện cơ chế biến hình(Form Switch). Khi nhấn phím chuyển dạng nhân vật chính có thể chuyển đổi thành hình dạng của quái vật bị đánh bại đồng thời sử dụng được hệ thống tấn công của quái vật ấy . Hình ~\ref{fig:bienhinh} hiển thị interact Prompt trên những xác quái vật có thể thực hiện cơ chế biến hình.
\begin{figure}[H]
    \centering
    \includegraphics[width=0.8\linewidth]{Hinhve/BienHinh.png}
    \caption{Interact Prompt hiện lên khi đứng gần xác quái vật để biến hình}
    \label{fig:bienhinh}
\end{figure}
Khi nhân vật trang bị các kỹ năng phép thuật, người chơi có thể thực hiện tiêu tốn mana để kích hoạt phép thuật hỗ trợ trong quá trình chiến đấu (Hình~\ref{fig:spellgameplay}).
\begin{figure}[H]
    \centering
    \includegraphics[width=0.8\linewidth]{Hinhve/SpellGameplay.png}
    \caption{Nhân vật sử dụng phép thuật trọng kiếm ánh sáng}
    \label{fig:spellgameplay}
\end{figure}

\subsection{Giai đoạn 5: Đối đầu với các boss cuối các khu vực}
Tại điểm cuối của các phân khu bản đồ lớn, người chơi bước vào các phòng Boss khổng lồ để chiến đấu. Các trận chiến đòi hỏi kỹ năng cao khi các Boss sở hữu lượng máu lớn cùng (Hiển thị ngay dưới màn hình một thanh HP riêng của Boss) và bộ não AI đặc thù để thực hiện các đòn tấn công đa dạng: đối đầu với Boss khổng lồ Big Golem tại khu vực Di tích (Hình~\ref{fig:bosstranchien}) và Boss Marionette tại khu vực Núi lửa (Hình~\ref{fig:bosstranchien2}).

\begin{figure}[H]
    \centering
    \includegraphics[width=0.8\linewidth]{Hinhve/Boss Golem.jpg}
    \caption{Đối đầu với Boss khổng lồ Big Golem}
    \label{fig:bosstranchien}
\end{figure}

\begin{figure}[H]
    \centering
    \includegraphics[width=0.8\linewidth]{Hinhve/BossMarionette.jpg}
    \caption{Đối đầu với Boss Marionette}
    \label{fig:bosstranchien2}
\end{figure}

\subsection{Các bảng giao diện chức năng quan trọng khác}
Để người đọc có thể hình dung toàn diện về các giao diện quan trọng khác trong "Dream Witch", các giao diện quản lý chỉ số nhân vật, cửa hàng mua sắm vật phẩm, cây kỹ năng và bản đồ định vị thế giới được trình bày tại các hình sau đây.

\begin{figure}[H]
    \centering
    \begin{subfigure}[b]{0.49\linewidth}
        \centering
        \includegraphics[width=\linewidth]{Hinhve/Status.jpg}
        \caption{Bảng trạng thái (Status Panel) hiển thị chỉ số chi tiết các linh hồn}
        \label{fig:statusui}
    \end{subfigure}
    \hfill
    \begin{subfigure}[b]{0.49\linewidth}
        \centering
        \includegraphics[width=\linewidth]{Hinhve/Shop.jpg}
        \caption{Giao diện Cửa hàng (Shopkeeper UI) mua sắm các phép thuật và công cụ}
        \label{fig:shopui}
    \end{subfigure}
    
    \vspace{10pt}
    
    \begin{subfigure}[b]{0.49\linewidth}
        \centering
        \includegraphics[width=\linewidth]{Hinhve/SkillTree.jpg}
        \caption{Nhánh kỹ năng (Skill Tree) mở khóa các đòn đánh nâng cao}
        \label{fig:skilltreeui}
    \end{subfigure}
    \hfill
    \begin{subfigure}[b]{0.49\linewidth}
        \centering
        \includegraphics[width=\linewidth]{Hinhve/FullMap.jpg}
        \caption{Bản đồ thế giới (Full Map UI) hỗ trợ người chơi định vị}
        \label{fig:fullmapui}
    \end{subfigure}
    \caption{Các giao diện quản lý thuộc tính, cửa hàng, cây kỹ năng và bản đồ thế giới}
    \label{fig:other_panels}
\end{figure}


\section{Kiểm thử}
\label{section:5.6}
Theo yêu cầu kiểm thử, nội dung phần này tập trung thiết kế các trường hợp kiểm thử (Test Cases) cho \textbf{ba chức năng quan trọng nhất} của trò chơi, bao gồm:
\begin{enumerate}
    \item Hệ thống điều khiển và di chuyển nhân vật (Player Movement).
    \item Cơ chế chuyển đổi hình thái linh hồn (Form Switch).
    \item Hệ thống lưu và tải tiến trình trò chơi (Save/Load System).
\end{enumerate}

\textbf{Kỹ thuật kiểm thử áp dụng:}
Quá trình kiểm thử áp dụng phương pháp \textbf{Kiểm thử hộp đen (Black-box Testing)}, sử dụng các kỹ thuật thiết kế ca kiểm thử cụ thể:
\begin{itemize}
    \item \textit{Phân vùng tương đương (Equivalence Partitioning):} Nhóm các trạng thái đầu vào (như trên mặt đất, trên không, sát tường) thành các vùng tương đương để thiết kế ca kiểm thử di chuyển, nhằm tối ưu hóa số lượng ca kiểm thử.
    \item \textit{Kiểm thử dựa trên kịch bản (Scenario Testing):} Giả lập luồng hành động thực tế của người chơi (Ví dụ: Tiêu diệt quái $\rightarrow$ Đứng gần xác $\rightarrow$ Biến hình) để kiểm tra tính toàn vẹn của chuỗi logic.
\end{itemize}

Chi tiết các trường hợp kiểm thử khác cho các hệ thống phụ (UI, Cửa hàng, Cây kỹ năng) nếu cần sẽ được đính kèm vào phần Phụ lục của tài liệu.

\subsection{Thiết kế các trường hợp kiểm thử cốt lõi}
Chi tiết 6 trường hợp kiểm thử đại diện cho ba chức năng quan trọng nhất được trình bày trong Bảng~\ref{table:test_cases}.

\begin{table}[H]
\centering
\begin{tabular}{|p{1.2cm}|p{2.5cm}|p{4.0cm}|p{4.5cm}|}
\hline
\textbf{Mã ca} & \textbf{Chức năng} & \textbf{Các bước thực hiện} & \textbf{Kết quả kỳ vọng} \\ \hline
\textbf{TC-01} & Di chuyển cơ bản & Nhấn giữ phím \texttt{A/D} trên bàn phím. & Nhân vật di chuyển mượt mà sang trái/phải, đổi hướng ngay lập tức và chạy animation chạy tương ứng. \\ \hline
\textbf{TC-02} & Nhảy đôi \& Leo tường & Nhấn \texttt{Space} 2 lần liên tiếp. Nhấn phím hướng ép vào tường. & Thực hiện nhảy đôi tạo lực đẩy lên trên. Nhân vật bám chặt vách tường và trượt xuống từ từ. \\ \hline
\textbf{TC-03} & Biến hình hợp lệ & Di chuyển tới xác kẻ địch đã chết, nhấn phím (\texttt{Tab}). & Đổi sang hình dáng quái vật, cập nhật đúng các thuộc tính (Max HP, tốc độ chạy, bộ kỹ năng mới). \\ \hline
\textbf{TC-04} & Hủy biến hình & Đang ở dạng quái vật, nhấn (\texttt{Tab}) lần nữa. & Hủy biến hình, nhân vật trở lại dạng người mặc định với thuộc tính ban đầu. \\ \hline
\textbf{TC-05} & Tự động lưu game & Điều khiển nhân vật đi qua cổng dịch chuyển (Portal) sang khu mới. & Hệ thống tự động ghi nhận trạng thái, xuất ra tệp JSON lưu trữ chính xác vị trí và trang bị. \\ \hline
\textbf{TC-06} & Tải lại tiến trình & Khởi chạy lại game, chọn ``Continue'' và chọn slot lưu tương ứng. & Đọc thành công tệp JSON, khôi phục nhân vật tại đúng tọa độ lưu gần nhất cùng toàn bộ vật phẩm. \\ \hline
\end{tabular}
\caption{Bảng đặc tả các trường hợp kiểm thử cho ba chức năng cốt lõi}
\label{table:test_cases}
\end{table}

\subsection{Tổng kết kết quả kiểm thử}
\label{subsection:5.6.1}
Trong đợt kiểm thử nội bộ (Alpha Test) đầu tiên, kết quả thực thi 6 ca kiểm thử trên nhiều cấu hình thiết bị khác nhau được tổng hợp chi tiết trong Bảng~\ref{table:test_results}.

\begin{table}[H]
\centering
\begin{tabular}{|p{3.5cm}|p{1.8cm}|p{1.8cm}|p{1.8cm}|p{2.6cm}|}
\hline
\textbf{Nhóm chức năng} & \textbf{Số ca} & \textbf{Đạt} & \textbf{Không đạt} & \textbf{Tỷ lệ thành công} \\ \hline
\textbf{Điều khiển \& Di chuyển} & 2 & 2 & 0 & 100\% \\ \hline
\textbf{Biến hình (Form Switch)} & 2 & 1 & 1 & 50\% \\ \hline
\textbf{Lưu/Tải (Save/Load)} & 2 & 2 & 0 & 100\% \\ \hline
\textbf{Tổng cộng} & \textbf{6} & \textbf{5} & \textbf{1} & \textbf{83.3\%} \\ \hline
\end{tabular}
\caption{Bảng tổng hợp kết quả kiểm thử hệ thống (Đợt 1)}
\label{table:test_results}
\end{table}

\textbf{Phân tích lý do kết quả kiểm thử không đạt:}
Tại lần kiểm thử đợt 1, ca kiểm thử \textbf{TC-03 (Biến hình hợp lệ)} cho ra kết quả \textbf{Không đạt}. Cụ thể:
\begin{itemize}
    \item \textbf{Hiện tượng lỗi (Bug):} Khi nhân vật người chơi đứng sát mép một bức tường hẹp hoặc trong góc kẹt và thực hiện biến hình thành một quái vật có kích thước lớn hơn (như Golem Boss), nhân vật bị kẹt vào bên trong khối va chạm (Collider) của bức tường. Sau đó nhân vật không thể di chuyển và bị vô hiệu hóa hoàn toàn.
    \item \textbf{Nguyên nhân gốc rễ:} Do Collider của Prefab quái vật mới lớn hơn nhiều so với kích thước Collider của nhân vật dạng người. Lệnh tải cấu trúc Prefab mới (\texttt{Instantiate}) diễn ra tức thời tại nguyên tọa độ gốc của nhân vật, dẫn đến vi phạm không gian vật lý của môi trường (Physics Overlap).
    \item \textbf{Khắc phục \& Đánh giá lại:} Nhóm phát triển đã sửa lỗi bằng cơ chế kiểm tra khoảng trống (Clearance Check). Trước khi chuyển đổi, hệ thống bắn tia \texttt{Raycast} xung quanh. Nếu phát hiện sát tường, hệ thống tự động đẩy nhẹ nhân vật ra xa một khoảng cách an toàn (Offset) rồi mới thay thế Collider. Sau khi áp dụng bản vá cập nhật mã nguồn, ca kiểm thử TC-03 đã hoạt động hoàn hảo và \textbf{đạt 100\% thành công} ở đợt kiểm thử hồi quy (Regression Testing).
\end{itemize}


\section{Triển khai}
\label{section:5.7}
Trò chơi được triển khai dưới dạng ứng dụng ngoại tuyến độc lập (PC Standalone) trên hệ điều hành Windows với các đặc tính kỹ thuật sau:

\begin{itemize}
    \item \textbf{Đóng gói sản phẩm (Build \& Packaging):} Sử dụng \textit{Unity Build Pipeline} để biên dịch sang mã máy Windows 64-bit, tạo ra tệp thực thi \texttt{DreamWitch.}\allowbreak\texttt{exe} kèm thư mục tài nguyên \texttt{DreamWitch\_Data}.
    \item \textbf{Phương thức phân phối (Distribution):} Sản phẩm được nén (\texttt{.zip}/\texttt{.rar}) để phân phối qua Google Drive hoặc Itch.io. Người chơi chỉ cần tải về, giải nén và vận hành trực tiếp tệp thực thi mà không cần cài đặt.
    \item \textbf{Cấu hình vận hành (Configuration):} Kích hoạt mặc định V-Sync giúp ổn định hiệu năng (khóa ở 60 FPS hoặc theo tần số quét màn hình). Độ phân giải tự động thích ứng theo tỷ lệ 16:9 (từ 720p đến 1080p) thông qua thành phần \texttt{Canvas Scaler}.
\end{itemize}

\end{document}
